\section{功能响应表}

本章对照课程实习技术评分表（50 分制，7 个评分维度）逐项说明技术响应情况。每条响应包含：需求项、响应情况（已实现/部分实现/待补充/后续规划）、实现说明和验证方式。

\begin{longtable}{|p{0.5cm}|p{3cm}|p{1.8cm}|p{5cm}|p{3.5cm}|}
  \caption{功能需求响应表}
  \label{tab:func-response}\\
  \toprule
  \textbf{\#} & \textbf{评价项} & \textbf{响应情况} & \textbf{实现说明} & \textbf{验证方式} \\
  \midrule
  \endfirsthead

  \multicolumn{5}{c}{\tablename~\thetable\quad（续）}\\
  \toprule
  \textbf{\#} & \textbf{评价项} & \textbf{响应情况} & \textbf{实现说明} & \textbf{验证方式} \\
  \midrule
  \endhead

  \bottomrule
  \multicolumn{5}{r}{\footnotesize 续下页}\\
  \endfoot

  \bottomrule
  \endlastfoot

  % ===== 维度1：项目定位与需求响应（6分） =====
  1.1 & 围绕五个核心用户问题说明用户场景 & 已实现 & 在产品计划书、功能设计书和本投标文件中明确阐述“去哪钓、能不能钓、现在适不适合钓、钓到了怎么分享和复盘、新手怎么快速上手”五个场景及对应承接模块 & 检查项目概述（第1章）和需求分析（第2章）中的用户场景描述 \\
  \hline
  1.2 & 明确 MVP 边界，不把非 MVP 内容作为当前交付 & 已实现 & 开发设计书和本投标文件明确列出 MVP 范围（地图、POI、记录、社区、推荐、教程）和非 MVP 范围（会员、赛事、商户、遥感、支付等），总体方案章节明确“当前 MVP 阶段已实现”与“后续可扩展”的区别 & 检查开发设计书 MVP 范围声明及本文件第1、3章的边界说明 \\
  \hline

  % ===== 维度2：核心功能完成度（15分） =====
  2.1 & 地图首页与 POI 展示 & 已实现 & 前端 HomeView.vue 与 MapPanel.vue 集成高德 JS API 2.0，展示附近 POI 标记、支持搜索、地图热力表达和天气摘要条；POI 数据通过后端 `/api/pois` 接口从 `data/pois.json` 加载 & 前端运行后进入首页地图，检查 POI 标记、搜索功能和天气摘要 \\
  \hline
  2.2 & 钓点详情 & 已实现 & 点击 POI 展示详情弹窗/页面：名称、类型、距离、鱼种、设施、评分分项、推荐理由、风险提示；数据来自 \texttt{/api/pois/\{poi\_id\}} 接口 & 检查 POI 详情字段完整性，对照 FishingPOI 数据模型 \\
  \hline
  2.3 & 鱼获/空军发布 & 已实现 & RecordView.vue 实现开竿计时和收竿汇总，收竿弹窗支持填写鱼种、条数、重量、钓法、饵料；条数为 0 时标记为空军；后端 `/api/records` POST 接口持久化到 `data/records.json` & 完成一次出钓记录（含鱼获和空军两种情况），重启后端检查数据是否仍存在 \\
  \hline
  2.4 & 社区鱼获流 & 已实现 & CommunityView.vue 展示社区帖子流，支持频道切换（推荐/同城/话题）；PublishView.vue 支持图文发布、关联钓点、内容与位置权限分离；后端 `/api/posts` 和 `/api/feed` 接口 & 发布一条帖子后在社区流中可见，检查帖子与出钓记录的关联 \\
  \hline
  2.5 & 教程页 & 已实现 & TutorialsView.vue 展示教程列表，支持分类筛选（入门/进阶、图文/视频）；数据来自 `data/tutorials.json` & 检查教程列表渲染、分类切换功能 \\
  \hline
  2.6 & 完整流程联动 & 部分实现 & “地图找点 → 详情 → 记录 → 从记录发布 → 社区展示”路径已跑通；记录到帖子的关联（record\_id）已在数据模型中预留并由发布弹层快速填充；个人统计已由前端展示结合后端 profile-summary/report 接口支撑 & 逐步骤演示完整流程，检查各页面间数据传递 \\
  \hline

  % ===== 维度3：技术架构与接口设计（8分） =====
  3.1 & 前后端分离架构 & 已实现 & Vue 3 + Vite 前端与 FastAPI 后端通过 RESTful API 通信；Vite 代理 `/api` 到后端 8000 端口；前端 API 调用集中在 `api.js` & 检查 `vite.config.js` 代理配置和后端路由注册 \\
  \hline
  3.2 & 后端三层架构 & 已实现 & 路由层（`app/api/`）→ 服务层（`app/services/`）→ 数据层（\texttt{app/data/json\_store.py}），职责清晰 & 检查后端目录结构和模块间调用关系 \\
  \hline
  3.3 & API 接口完整性 & 已实现 & 已实现 POI、帖子、记录、推荐、天气、AI 解释、教程等核心接口，统一使用 `data` + `meta` 响应格式 & 通过浏览器开发者工具或接口测试工具检查各端点返回结构 \\
  \hline
  3.4 & 环境变量与密钥管理 & 已实现 & 高德 Key、DeepSeek Key 通过 `.env` 文件由 `config.py` 读取，不写入前端代码；`.env.example` 提供模板 & 检查 `.env` 不在版本控制中，前端无硬编码 Key \\
  \hline

  % ===== 维度4：推荐模型与AI解释设计（7分） =====
  4.1 & 规则模型排序 & 已实现 & 推荐接口 `/api/recommendations` 实现六维评分（鱼情分 0.25、合规分 0.25、距离分 0.15、天气适配分 0.15、设施便利分 0.10、社交热度分 0.10），包含硬性拦截/惩罚项（禁钓、极端天气、岸线安全等） & 检查推荐接口返回的 score、breakdown 和 warnings 字段 \\
  \hline
  4.2 & AI 解释设计 & 部分实现 & AI 接口 `/api/ai/fishing-advice` 已预留结构：接收候选 POI、规则推荐结果和用户意图等结构化数据，返回 JSON 格式的自然语言解释；当前 MVP 阶段若在 production 环境配置 `DEEPSEEK\_API\_KEY` 则可真实调用，否则返回 mock 解释；AI 不参与排序，仅解释推荐结果 & 检查 AI 接口输入参数（不包含用户隐私原文）和输出格式（summary + explanation + top\_recommendations + tips） \\
  \hline
  4.3 & 风险优先提示 & 已实现 & 推荐结果中硬性风险会通过 penalty、warnings 和 POI 内的 risk/risk\_flags 字段体现；存在禁钓/限钓/安全风险的点位会降低排序并提示用户先核对规则与现场安全 & 检查包含禁钓 POI 时的推荐返回结果 \\
  \hline

  % ===== 维度5：数据结构与数据可用性（5分） =====
  5.1 & 核心数据模型 & 已实现 & FishingPOI（20+ 字段，含评分分项和风险标记）、FishingRecord（鱼种/重量/钓法/天气/隐私级别）、CatchPost（帖子内容、关联记录、内容与位置权限）、Recommendation（排名/分项分/理由/警告）、WeatherSnapshot 等模型通过 Pydantic 定义，字段能支撑核心流程 & 检查 `app/schemas/` 目录和接口返回字段 \\
  \hline
  5.2 & JSON 集合与 Schema & 已实现 & 根目录 `data/` 下每个 `.json` 集合配有同名 `.schema.json` 描述文件；当前包含 pois、posts、records、tutorials、weather\_snapshots 等集合 & 检查 `data/` 目录和 schema 文件内容 \\
  \hline
  5.3 & 数据持久化 & 已实现 & \texttt{json\_store.py} 的 \texttt{save\_collection()} 采用临时文件 + 原子替换策略，避免写入中断导致数据损坏；记录和帖子的增删改操作均调用持久化 & 新增记录后重启后端，检查数据仍存在且不重复 \\
  \hline

  % ===== 维度6：界面交互与演示体验（5分） =====
  6.1 & 页面结构与跳转 & 已实现 & 5 个底部 Tab（首页/社区/记录/技巧/服务）覆盖主要页面；“我的”通过顶部菜单抽屉进入，发布通过社区页弹层进入；地图即首页策略，Tab 切换无路由丢失 & 逐 Tab 切换并打开发布弹层、我的抽屉，检查各页面可正常展示 \\
  \hline
  6.2 & 状态反馈 & 部分实现 & 部分组件有加载和空状态处理；错误兜底和离线状态提示尚不完善，当前阶段主要依赖 mock 数据保证演示链路不断裂 & 断网或接口异常时检查页面是否展示错误提示而非白屏 \\
  \hline
  6.3 & 视觉统一性 & 已实现 & 统一使用 Element Plus 组件库，配色和布局保持一致性 & 检查各页面截图，对照视觉风格 \\
  \hline

  % ===== 维度7：文档规范与团队协作（4分） =====
  7.1 & 文档完整性 & 已实现 & 已输出产品计划书、功能设计书、开发设计书、技术评分表、技术投标文件（本文件），覆盖课程要求文档清单 & 对照课程文档要求逐项核对 \\
  \hline
  7.2 & 内容与实现一致 & 已实现 & 文档描述的功能模块、数据模型、API 接口、页面结构与实际代码一致；非 MVP 内容明确标注为“后续规划” & 随机抽取文档中的功能描述，在代码中验证是否存在对应实现 \\
  \hline
  7.3 & 编号与格式统一 & 已实现 & 本文件采用连续章节编号，表格使用统一样式（tabularx/longtable），功能响应表包含“需求项、响应情况、实现说明、验证方式”四列 & 检查文件目录、交叉引用和表格格式 \\
  \bottomrule
\end{longtable}
